% Section 6. Budget 3.75 pages.
% Source: draft/bie-enclosure/s-results.tex (clean version); data from
% test/bie-bound/pipeline-v1 run-002, resolution-v1, example2-v1 run-015,
% gapedge-v1 run-006/007/011/012, test/fem-goerisch/pipeline-v1.
\section{Numerical results}\label{sec:results}
Each interval below comes from one run that starts from an empty output
directory and performs, in order, the gap classification, the determinant
search, the density reconstruction and field fit, the reliable evaluation of
all five groups, the mass bound, the trace constants, the assembly, and the
gap test, without reading earlier results. Geometry parameters that are not
binary64 numbers (such as $0.2$) are replaced by the nearest binary64 number,
and the intervals refer to this geometry. Timings are wall-clock on an Apple M3
Pro (11 cores, 18\,GiB), MATLAB R2023b and python-flint 0.9.0. The code, the
saved candidate data and an independent checker are in the release package
\todo{cite release}.

\begin{figure}[tbhp]
\centering
\begin{tikzpicture}
\begin{groupplot}[group style={group size=1 by 2, vertical sep=1.05cm},
  width=\textwidth, height=0.23\textwidth, axis equal image,
  ymin=-0.5, ymax=0.5, ytick={-0.5,0,0.5}, xtick={-3,...,3},
  ylabel={$y$}, unbounded coords=jump, tick label style={font=\scriptsize},
  label style={font=\scriptsize}, title style={font=\scriptsize}]
\nextgroupplot[xmin=-3.5, xmax=3.5, title={Example 1}]
  \foreach \w in {-2.5,-1.5,...,2.5} {
    \edef\temp{\noexpand\draw[gray, dashed] (axis cs:\w,-0.5) -- (axis cs:\w,0.5);}
    \temp}
  \pgfplotsinvokeforeach{-3,-2,-1}{\addplot[blue, thick, fill=blue!15] table[col sep=comma,
    x=x, y expr={\thisrow{cell}==#1 ? \thisrow{y} : nan}] {data/geometry.csv};}
  \addplot[red, thick, fill=red!15] table[col sep=comma,
    x=x, y expr={\thisrow{cell}==0 ? \thisrow{y} : nan}] {data/geometry.csv};
  \pgfplotsinvokeforeach{1,2,3}{\addplot[teal, thick, fill=teal!15] table[col sep=comma,
    x=x, y expr={\thisrow{cell}==#1 ? \thisrow{y} : nan}] {data/geometry.csv};}
\nextgroupplot[xmin=-3.5, xmax=3.5, xlabel={$x$}, title={Example 2 ($L_0=1$)}]
  \foreach \w in {-2.5,-1.5,-0.5,0.5,1.35,2.2,3.05} {
    \edef\temp{\noexpand\draw[gray, dashed] (axis cs:\w,-0.5) -- (axis cs:\w,0.5);}
    \temp}
  \pgfplotsinvokeforeach{-3,-2,-1}{\addplot[blue, thick, fill=blue!15] table[col sep=comma,
    x=x, y expr={\thisrow{cell}==#1 ? \thisrow{y} : nan}] {data/geometry2.csv};}
  \pgfplotsinvokeforeach{1,2,3}{\addplot[teal, thick, fill=teal!15] table[col sep=comma,
    x=x, y expr={\thisrow{cell}==#1 ? \thisrow{y} : nan}] {data/geometry2.csv};}
  % Fourth right rod: the third one shifted by one period L_+ = 0.85 (clipped).
  \addplot[teal, thick, fill=teal!15] table[col sep=comma,
    x expr={\thisrow{x}+0.85}, y expr={\thisrow{cell}==3 ? \thisrow{y} : nan}] {data/geometry2.csv};
\end{groupplot}
\end{tikzpicture}
\caption{Part of the strip in the two examples, on the same $x$ scale. Dashed
lines are the artificial walls; $y=\pm1/2$ is the seam. The left leads are the
same; in Example~2 the center cell is empty and the right lead has period
$L_+=0.85$ and a rotated rod. Lead rods have $\rho=17$, the star rod $\rho=4$.}
\label{fig:geometry}
\end{figure}

\subsection{Example 1: a star-shaped defect}\label{sec:ex1}
Let $d=1$, $\beta=0.5$, $L_0=L_\pm=1$, $\rho_{\mathrm e}=1$. The center cell
contains a star with $\rho=4$, and every lead cell a three-lobed rod with
$\rho=17$, centered in the cell,
\[
  \boldsymbol r_0(t)=0.2\begin{pmatrix}(1.5+0.3\cos5t)\cos t\\(1+0.3\cos5t)\sin t\end{pmatrix},
  \qquad
  \boldsymbol r_a(t)=\frac{0.2(1+a\cos3t)}{(1+a^2/2)^{1/2}}\begin{pmatrix}\cos t\\\sin t\end{pmatrix},
\]
with $a=0.12$ on the left and $a=0.20$ on the right (\cref{fig:geometry}). Thus $\rin=1$, $\rax=17$ and
$\lambda_0=1/68$.

\emph{Discretization.} The center curve has $512$ nodes, each lead curve $256$;
$M=48$ ($K=97$). The trial field uses $480$ proxies (\cref{sec:field}). The
reliable evaluation splits the center curve into $192$ panels, each lead curve
into $64$, each wall into $16$ and each seam into $8$, with $48$ Chebyshev nodes per
panel and $128$-bit balls. Both leads use $p=1$, with
$\|\mathsf P_\pm\|_F\le0.3815$, $0.3868$; $N_{\mathrm{cell}}$ is the smallest depth
at which the squared tail is at most $10^{-6}$ times the squared finite part
($52$--$58$). The mass bound uses $8$ disks inside the rods. The gap test uses
$169$ plane waves and $118$ phase panels on $[0,\pi]$; both lead rods are
symmetric under $y\mapsto-y$.

\emph{Candidate.} The floating-point gap classification gives the search window
$k\in[1.451,2.249]$. Chebyshev interpolation of $\det\mathsf A(k)$ on the whole
window stops at degree $32$ (last three coefficients below $7\times10^{-11}$
relative to the largest), with one near-real root in the Bernstein ellipse. Two
Newton steps give $\widehat k=1.558437894182088=438661269969907/2^{48}$; $\mu$ is
its exact square. The search takes $38$ matrix evaluations and $159$\,s.

\emph{Interval.} \Cref{tab:ex1} lists the five groups. The center rod
contributes $90\%$ of $Q$. There the collar is narrow ($h_g=0.969$) because the
concave points of the star have curvature about $64$; this gives the largest
$C_g$. The interpolation errors satisfy $e_I\le2.2\times10^{-17}$ on all panels, so
the reliable $\overline\varepsilon_g$ exceed the floating-point values by
only $0.04\%$. Increasing every $N_{\mathrm{cell}}$ by $10$ changes the group bounds by a
relative $4.7\times10^{-7}$ at most. The disks give $\underline U=1.23587$, and
\begin{align*}
  R_{\mathrm{spec}}&\le1.3690\times10^{-11},\\
  I_{\mathrm{spec}}&\subset[2.42872867000901,\,2.42872867003640],\\
  I_k&\subset[1.55843789417769,\,1.55843789418648],
\end{align*}
with outward rounding. All $472$ inertia tests pass, so neither lead has spectrum
for $k\in(1.318,2.375)$; the margins in \eqref{eq:gap-condition} are $0.6916$ and
$3.2118$. Hence $I_k$ contains at least one guided-mode wavenumber. The complete
run takes $827$\,s with a peak memory of $1.84$\,GiB; $528$\,s of this is the
reliable evaluation on $8$ processes (\cref{sm:results}).

\begin{table}[tbhp]
\caption{Example 1: trace constants and defect bounds of the five groups (upward
rounded). Lead and wall groups include all cells of both half guides.}
\label{tab:ex1}
\centering\small
\begin{tabular}{lccc}
\toprule
Group $g$ & $\overline C_g$ & $\overline\varepsilon_g$ & $\overline C_g\overline\varepsilon_g$\\
\midrule
center rod & $2.490254$ & $1.2616\times10^{-12}$ & $3.1417\times10^{-12}$\\
left rods & $1.181583$ & $5.0811\times10^{-14}$ & $6.0037\times10^{-14}$\\
right rods & $1.233941$ & $5.2316\times10^{-14}$ & $6.4555\times10^{-14}$\\
walls & $1.035553$ & $2.0346\times10^{-13}$ & $2.1070\times10^{-13}$\\
seams & $1.035553$ & $5.7837\times10^{-15}$ & $5.9894\times10^{-15}$\\
\midrule
$Q$ & & & $3.4830\times10^{-12}$\\
\bottomrule
\end{tabular}
\end{table}

\emph{Resolution.} We vary the field and the evaluation separately (complete
reruns, \cref{sm:results}). With $256$, $384$, $512$ center nodes (half as many per
lead curve) the full width $2R_{\mathrm{spec}}$ is $4.28\times10^{-6}$,
$9.87\times10^{-9}$, $2.74\times10^{-11}$, almost entirely through the center group;
$M=40$ or $56$ changes the width by about $1\%$. The five candidates differ from
each other by at most five ulps. For the same field, doubling all panels,
$64$ Chebyshev nodes or $192$-bit balls change each group bound by a relative
$3\times10^{-6}$ at most. The width therefore measures how well the reconstructed
field satisfies the interface conditions, not the accuracy of $\widehat k$ or of
the evaluation.

\emph{Comparison with a finite element enclosure.} The same problem was solved
independently with $P_3$ finite elements on a mesh with $h=1/48$ (about
$2.9\times10^4$ degrees of freedom per cell), a trial field built from a discrete
DtN map and tapered to zero over $13$ lead cells, and Goerisch's method
\todo{cite FEM route}. Under a spectral isolation condition, checked by separate
finite-dimensional tests, it gives $\mu\in[2.4287286691977,2.4287286736602]$ (width
$4.46\times10^{-9}$) together with uniqueness in this interval; truncation errors
are bounded reliably but rounding is only estimated. $I_{\mathrm{spec}}$ lies
inside. The two codes define the binary geometry slightly differently (by
$2^{-56}$ in one Fourier coefficient), so this is a numerical cross-check rather
than a statement about one operator. The FEM run takes $12.5$\,min plus
$54.3$\,min for the isolation condition.
\todo{FEM cost--width series over three mesh sizes (release R3).}

\subsection{Example 2 and guided modes near a gap edge}\label{sec:ex2}
The second example removes the center rod, uses different lead periods, and
breaks the reflection symmetry. The center cell $|x|<L_0/2$ is empty; the left
lead has $L_-=1$ and the rod $a=0.12$; the right lead has $L_+=0.85$ and the rod
$a=0.20$ rotated by $0.13$ (\cref{fig:geometry}). The right bands at $\pm\theta$
then differ by up to $3.7\times10^{-3}$, and the gap test covers the full range
$[-\pi/L_s,\pi/L_s]$ with $236$ panels per lead. The discretization is the same
as in Example~1. The center field is a finite Rayleigh sum (\cref{thm:center-rep}
with $\Omega_0=\emptyset$), which is exactly quasi-periodic, so the center rod and
center seam groups vanish identically, and the walls dominate $Q$ ($57\%$). The
floating-point common gap is about $k\in(1.441,2.396)$.

Increasing $L_0$ lowers the guided wavenumber toward the lower gap edge, and the
decay in the right lead slows down. \Cref{tab:ex2} reports complete reliable
runs for $L_0=1,1.6,1.9,2$; at each point the lowest confirmed root in the window
is used, not continued from the previous one. Along a floating-point continuation,
$\ln(1/\varrho(\mathsf P_+))\approx(1.94\text{--}2.04)\sqrt{\delta_{\mathrm e}}$, where $\varrho$ is the
spectral radius and $\delta_{\mathrm e}$ the distance to the band edge $1.441$, in line with the square-root behavior of
evanescent decay near a quadratic band edge. Correspondingly
$N_{\mathrm{cell}}\ln(1/\overline r_+)\approx51,49,48$, so the depth grows like
$\delta_{\mathrm e}^{-1/2}$. The width stays below $10^{-11}$
because the tail enters only through its bound. From $L_0=1$ to $2$ the right
depth grows fourfold but the total time only by $42\%$, since the search and the
gap test do not depend on $L_0$. Throughout,
$\|\mathsf P_\pm\|_2/\varrho(\mathsf P_\pm)\le1.021$ and $p=1$ suffices; the $p=2$
bound \eqref{eq:tail}, evaluated on the same states, differs by about $2\%$. For
$L_0=2$ the gap test window is $(1.45,2.37)$, the smallest inertia margin is
$0.197$, and the interval has margin $0.0161$ to the lower end; a second
candidate $\widehat k\approx2.2156$ appears in the window and is not treated.

\begin{table}[tbhp]
\caption{Example 2 for increasing center width $L_0$: interval $I_k$ (shared
leading digits, then outward-rounded last digits), distance $\delta_{\mathrm e}$ of
$\widehat k$ to the floating-point band edge, upper bound $\overline r_+$ of
$\|\mathsf P_+\|_2$, right depths $N_{\mathrm{cell}}$ (rods/walls/seams), $Q$, full
width of $I_{\mathrm{spec}}$ (bounds rounded upward), and time of one cold run of the release package. The depths are chosen so that the
squared tail is at most $\tau$ times the squared finite part, with
$\tau=10^{-6},10^{-2},0.1,0.1$; the tail always enters through its bound.}
\label{tab:ex2}
\centering\footnotesize\setlength{\tabcolsep}{3pt}
\begin{tabular}{cccccccc}
\toprule
$L_0$ & $I_k$ & $\delta_{\mathrm e}$ & $\overline r_+$ & $N_{\mathrm{cell}}$ & $Q$ & width & s\\
\midrule
$1$   & $1.89169614005[23,32]$ & $0.451$  & $0.3344$ & $52/51/52$ & $1.40\times10^{-13}$ & $3.09\times10^{-12}$ & $1076$\\
$1.6$ & $1.58144476391[49,70]$ & $0.140$  & $0.4924$ & $72/69/72$ & $4.13\times10^{-13}$ & $6.36\times10^{-12}$ & $1194$\\
$1.9$ & $1.48029280022[43,61]$ & $0.0393$ & $0.6825$ & $128/121/126$ & $3.87\times10^{-13}$ & $5.22\times10^{-12}$ & $1353$\\
$2$   & $1.45555716328[31,60]$ & $0.0146$ & $0.7953$ & $208/196/206$ & $6.24\times10^{-13}$ & $8.16\times10^{-12}$ & $1532$\\
\bottomrule
\end{tabular}

\end{table}
